\documentclass[11pt, a4paper]{article}

% Gemeinsame Style-Definitionen (TEXINPUTS muss templates/ enthalten — siehe scripts/build_tex.py)
\input{bfw-style.tex}

\title{\vspace*{-1.5cm}\Huge\bfseries\color{bfwPrimaryDark}Session~5: Geschäftsprozesse, Produktionsfaktoren \& Güterarten\\[0.4em]
  \large\color{bfwInkSoft}\textnormal{Wertschöpfungskette, Prozessarten und Güterklassifikation bei JIKU IT-Solutions}}
\author{\textsf{IT Lernfeld 1 --- Das Unternehmen und die eigene Rolle im Betrieb beschreiben}}
\date{}

\begin{document}
\maketitle
\thispagestyle{empty}

\vspace{1em}
\begin{merksatz}[title={\bfseries Worum geht es in dieser Session?}]
Jedes Unternehmen erzeugt \emph{Wertschöpfung} --- den Mehrwert, der nach Abzug aller
eingekauften Vorleistungen und Abschreibungen verbleibt. Dieser Wert kommt Mitarbeitern,
Fremdkapitalgebern, dem Staat und den Eigentümern zugute. Am Beispiel von \emph{JIKU
IT-Solutions} lernen wir, wie Geschäftsprozesse in Kern-, Support- und Führungsprozesse
eingeteilt werden, welche \emph{betriebswirtschaftlichen Produktionsfaktoren} nach Gutenberg
unterschieden werden und wie sich Güter nach verschiedenen Kriterien klassifizieren lassen.
\end{merksatz}

\tableofcontents
\vspace{1em}
\hrule
\vspace{1em}

\section{Wertschöpfung und Wertschöpfungskette}

\subsection{Begriff und Berechnung der Wertschöpfung}

Ziel jedes Unternehmens ist es, zwischen Input und Output eine möglichst große Wertschöpfung
zu erzielen. Wertschöpfungsintensiv sind vor allem die Tätigkeiten, die direkt mit der
Leistungserstellung zu tun haben, kundennah sind und zu den Kernkompetenzen des Unternehmens
gehören.

\begin{center}
\begin{tabular}{p{4cm} p{9cm}}
\toprule
\textbf{Begriff} & \textbf{Erläuterung}\\
\midrule
Wertschöpfung & Umsatzerlöse (netto) +/- Lagerveränderung abzüglich Vorleistungen abzüglich Abschreibungen\\
\addlinespace
Vorleistungen & Einkauf jeglicher Vorprodukte oder Dienstleistungen für den betrieblichen Zweck; ohne Löhne, Zinsen und Steuern\\
\addlinespace
Abschreibungen & Wertminderungen auf Anlage- und Umlaufvermögen (Gebäude, Maschinen, Systeme, Fahrzeuge, Stoffe, Teile, Handelswaren)\\
\addlinespace
Wertschöpfungstiefe & Wertschöpfung geteilt durch Gesamtleistung (Output); zeigt, wie viel des Umsatzes das Unternehmen selbst erwirtschaftet\\
\addlinespace
W/M-Kennzahl & Wertschöpfung pro Mitarbeiter = Wertschöpfung geteilt durch Mitarbeiterzahl; Maß für Produktivität\\
\bottomrule
\end{tabular}
\end{center}

\begin{merksatz}[title={\bfseries Merke: Wertschöpfungsnutznießer}]
Die Wertschöpfung kommt vier Gruppen zugute:
\begin{enumerate}
  \item \textbf{Mitarbeiter} --- Löhne, Gehälter, Sozialabgaben, Leistungsprämien
  \item \textbf{Fremdkapitalgeber} --- Zinsen für Darlehen
  \item \textbf{Staat} --- Steuern und Abgaben
  \item \textbf{Eigenkapitalgeber} --- Gewinne als Risikoprämie und Investitionskapital
\end{enumerate}
Sinkt die Wertschöpfung, können diese Stakeholder nicht mehr mit den bisherigen Beteiligungen
rechnen --- es kommt zu Verteilungsauseinandersetzungen.
\end{merksatz}

\subsection{Wertschöpfungskette nach Branchen}

Die Wertschöpfungskette zeigt die wertschöpfungsintensiven Tätigkeitsbereiche je Branche.
Verwaltende Tätigkeiten außerhalb der Kette (z.\,B. Personal, Rechnungswesen, interne IT)
zählen \emph{nicht} zur Wertschöpfungskette.

\begin{center}
\begin{tabular}{p{3.5cm} p{9.5cm}}
\toprule
\textbf{Betriebstyp} & \textbf{Wertschöpfungskette}\\
\midrule
Industriebetrieb & Beschaffung → Herstellung → Marketing → Verkauf → Service\\
\addlinespace
Handelsbetrieb & Beschaffung → Lagerung → Marketing → Verkauf → Service\\
\addlinespace
Dienstleistungs\-betrieb & Marketing → Verkauf → Beschaffung → Dienstleistung → Service\\
\bottomrule
\end{tabular}
\end{center}

\subsection{Digitalisierung und Wertschöpfung}

\begin{center}
\begin{tabular}{p{3.5cm} p{9.5cm}}
\toprule
\textbf{Begriff} & \textbf{Definition}\\
\midrule
Digitalisierung & Veränderungsprozess, bei dem bisher analoge Prozesse und Daten künftig digital bearbeitet werden.\\
\addlinespace
Digitale Transformation & Veränderte Wertschöpfungs- und Geschäftsprozesse, die durch digitale Technologien oder entsprechende Kundenerwartungen entstanden sind.\\
\addlinespace
Digitale Disruption & Löst ein bestehendes Kundenproblem auf eine komplett neue Art und schafft damit ganz andere Geschäftsmodelle; geht weiter als Transformation.\\
\bottomrule
\end{tabular}
\end{center}

\section{Geschäftsprozesse im Betrieb}

\subsection{Prozessarten}

Im Unternehmen lassen sich alle Aktivitäten in Geschäftsprozesse einteilen. Diese bestehen
aus Sub- und Unterprozessen und auf der untersten Detaillierungsebene aus Workflows.

\begin{center}
\begin{tabular}{p{3.5cm} p{9.5cm}}
\toprule
\textbf{Prozessart} & \textbf{Definition und Beispiele}\\
\midrule
Geschäftsprozess & Zeitliche und sachlogische Abfolge von Unternehmensaktivitäten, die festgelegte Unternehmensziele verfolgen und auf Unternehmensressourcen zurückgreifen.\\
\addlinespace
Kernprozess (Wertschöpfungs\-prozess) & Kundennah und besonders wertschöpfungsintensiv; das Unternehmen weist dafür besondere Kernkompetenzen nach. Beispiele JIKU: Aufbau komplexer IT-Systeme, IT-Serviceleistungen, Auftragsabwicklung.\\
\addlinespace
Supportprozess (Unterstützungs\-prozess) & Nicht direkt wertschöpfend, aber notwendig, um Kernprozesse zu ermöglichen. Beispiele: Buchhaltung, Lohn- und Gehaltsabrechnung, Controlling.\\
\addlinespace
Workflow & Teilprozess auf der untersten Detaillierungsebene; läuft rechnergestützt und weitgehend automatisiert ab; gesteuert durch ein Workflowmanagementsystem.\\
\bottomrule
\end{tabular}
\end{center}

\begin{merksatz}[title={\bfseries JIKU-Praxisbeispiel: Leistungs-, Geld- und Informationsflüsse}]
Im Schaubild der Leistungs-, Geld- und Informationsflüsse bei JIKU lassen sich alle
Einzelprozesse zuordnen: Lieferanten beliefern das Lager und die Werkstatt; Einkauf und
Verkauf koordinieren die Warenflüsse; Rechnungswesen und Controlling steuern Geldflüsse.
Viele dieser Prozesse laufen heute digital ab.
\end{merksatz}

\subsection{Integrierte Unternehmenssoftware}

Entlang der Wertschöpfungskette werden integrierte Softwaresysteme eingesetzt.

\begin{center}
\begin{tabular}{p{1.5cm} p{11.5cm}}
\toprule
\textbf{Kürzel} & \textbf{Bedeutung}\\
\midrule
E-P & E-Procurement: elektronische Beschaffung über Internet und Beschaffungsplattformen\\
SRM & Supplier-Relationship-Management: Lieferantenbeziehungsmanagement\\
SCM & Supply-Chain-Management: Verwaltung integrierter Logistikketten über den gesamten Wertschöpfungsprozess\\
PPS & Produktionsplanung und -steuerung\\
CRM & Customer-Relationship-Management: Ausrichtung auf Kunden und Gestaltung von Kundenbeziehungsprozessen\\
ERP & Enterprise-Resource-Planning: Verwaltung der Unternehmensressourcen (Lager, Auftragsbearbeitung, Gehaltsabrechnung, Rechnungswesen)\\
DMS & Dokumentenmanagement: datenbankgestützte Verwaltung elektronischer Dokumente\\
BPM & Business-Process-Management: Steuerung und Optimierung von Geschäftsprozessen\\
BI & Business Intelligence: entscheidungsrelevante Unternehmensdaten für automatisierte Reports\\
\bottomrule
\end{tabular}
\end{center}

\section{Betriebliche Produktionsfaktoren}

\subsection{Volkswirtschaftliche vs. betriebswirtschaftliche Betrachtung}

Unter Produktionsfaktoren versteht man alle materiellen und immateriellen Mittel und Leistungen,
die an der Bereitstellung von Gütern mitwirken. Es wird zwischen volkswirtschaftlicher
(Arbeit, Boden, Kapital) und betriebswirtschaftlicher Betrachtung unterschieden.

\subsection{Das Gutenberg-Modell}

Der Wirtschaftswissenschaftler \textbf{Erich Gutenberg} (1897--1984) unterschied für
Produktionsbetriebe folgende Faktoren:

\begin{center}
\begin{tabular}{p{3.5cm} p{9.5cm}}
\toprule
\textbf{Faktor} & \textbf{Definition und Beispiele}\\
\midrule
Werkstoffe & Alle Güter, die zur Erstellung eines Produktes notwendig sind: Rohstoffe (Hauptbestandteile), Hilfsstoffe (Nebenbestandteile, z.\,B. Klebstoff), Betriebsstoffe (nicht direkt zurechenbar, z.\,B. Schmieröl), Fertigteile\\
\addlinespace
Betriebsmittel & Technische Ausstattungen: Betriebsgebäude, Maschinen, Systeme, Fahrzeuge\\
\addlinespace
Ausführende Arbeit & Menschliche Arbeitsleistung \emph{ohne} Leitungs-, Planungs-, Organisations- oder Kontrollfunktion\\
\addlinespace
Dispositiver Faktor & Die Geschäftsleitung: ergänzt die Elementarfaktoren zu einer produktiven Einheit; Aufgaben: Leitung, Planung, Organisation, Kontrolle\\
\addlinespace
Information (ergänzend) & Wissen, Rechte, Daten; wurde nach Gutenberg als eigenständiger Faktor ergänzt, da die Informationswirtschaft seitdem stark gewachsen ist\\
\bottomrule
\end{tabular}
\end{center}

\begin{merksatz}[title={\bfseries Merke: Grenzen des Gutenberg-Modells}]
Gutenberg richtete sein Modell auf Produktionsbetriebe aus --- Handelsbetriebe und die
Informationswirtschaft blieben außen vor. Der Faktor \textbf{Zeit} (z.\,B. Lieferzeit im
Handel) und der Faktor \textbf{Information} (Know-how, Softwarelizenzen, Patente) wurden
erst später ergänzt. Für JIKU IT-Solutions ist Information der entscheidende Faktor.
\end{merksatz}

\subsection{Ökonomische Prinzipien}

\begin{center}
\begin{tabular}{p{3cm} p{5.5cm} p{4cm}}
\toprule
\textbf{Prinzip} & \textbf{Ziel} & \textbf{Beispiel}\\
\midrule
Maximumprinzip & Mit gegebenen Mitteln maximalen Erfolg erzielen & Mit 50~€ Budget möglichst viele Kunden erreichen\\
\addlinespace
Minimumprinzip & Bestimmten Erfolg mit minimalen Mitteln erzielen & Werbebotschaft mit geringstem Aufwand übermitteln\\
\addlinespace
Optimumprinzip (Extremum\-prinzip) & Bei variablem Ergebnis und Einsatz die optimale Kombination bestimmen & Angebot mit dem niedrigsten Wert Output/Input wählen\\
\bottomrule
\end{tabular}
\end{center}

\section{Güterarten}

\subsection{Klassifikationsmerkmale}

Güter lassen sich nach verschiedenen Kriterien einteilen:

\begin{center}
\begin{tabular}{p{3cm} p{4.5cm} p{5.5cm}}
\toprule
\textbf{Kriterium} & \textbf{Güterarten} & \textbf{Beispiele}\\
\midrule
Beschaffenheit & Materiell / Immateriell & Auto, Gebäude / Dienstleistung, Softwarelizenz\\
\addlinespace
Ort und Zweck & Konsumgut / Investitionsgut & Druckerpapier für Privatnutzung / Serverrack\\
\addlinespace
Verbrauch & Verbrauchsgut / Gebrauchsgut & Strom, Telefonminuten / Maschine, Regal\\
\addlinespace
Knappheit & Freies Gut / Wirtschaftliches Gut & Atemluft, Sonnenschein / Buch, Sauerstoff\-flasche\\
\addlinespace
Ausschließbarkeit & Privates Gut / Öffentliches Gut & Kinovorstellung / Parkanlage, Leuchtturm\\
\addlinespace
Beziehung & Komplementär / Substitutiv / Unverbunden & Drucker und Toner / Zug und Auto / PC und Blumen\\
\bottomrule
\end{tabular}
\end{center}

\subsection{Güter zur Leistungserstellung (Einkaufsperspektive)}

\begin{center}
\begin{tabular}{p{3cm} p{10cm}}
\toprule
\textbf{Kategorie} & \textbf{Erläuterung}\\
\midrule
Rohstoffe & Hauptbestandteile eines Erzeugnisses, z.\,B. Eisen für Gehäuse\\
Hilfsstoffe & Nebenbestandteile, z.\,B. Farbe, Klebstoff, Heftklammern\\
Betriebsstoffe & Nicht direkt einem Erzeugnis zurechenbar, z.\,B. Schmieröle, Brennstoffe\\
Einbauteile & Komponenten für die Herstellung eines Produktes\\
Handelswaren & Güter, die an den Endverbraucher weiterverkauft werden\\
Dienstleistungen & Fremdleistungen (Dienste), die zur Leistungserstellung notwendig sind\\
Lizenzen, Rechte & Markenrechte, Softwarenutzungsrechte, Patente und Urheberrechte\\
\bottomrule
\end{tabular}
\end{center}

\section{Organisationsmittel}

Unter Organisationsmitteln werden in der Aufbau- und Ablauforganisation sowohl
\textbf{Arbeitsmittel} (z.\,B. Formulare, Checklisten, Organigramme) als auch
\textbf{formale Gestaltungsmittel} (z.\,B. Stellenbeschreibungen, Prozesshandbücher)
verstanden. Mit dem Ziel der Digitalisierung stehen hier große Veränderungen an:
ERP-Systeme, Ticketsysteme, digitale Dokumentenmanagementsysteme und Workflowtools
ersetzen zunehmend papierbasierte Organisationsmittel.

\begin{merksatz}[title={\bfseries Zusammenfassung: Kernbegriffe dieser Session}]
\textbf{Wertschöpfung} = Output minus Vorleistungen minus Abschreibungen ~·~
\textbf{Kernprozesse} stiften Kundennutzen ~·~
\textbf{Supportprozesse} ermöglichen die Kernprozesse ~·~
\textbf{Elementarfaktoren} nach Gutenberg: Werkstoffe, Betriebsmittel, ausführende Arbeit ~·~
\textbf{Dispositiver Faktor}: Geschäftsleitung ~·~
\textbf{Güterklassifikation}: stets nach mehreren Kriterien gleichzeitig prüfen.
\end{merksatz}

\end{document}
